\documentclass[acmsmall,11pt,nonacm]{acmart}
\geometry{letterpaper,margin=1in,headheight=14pt,footskip=18pt}
\settopmatter{printacmref=false,printccs=false,printfolios=true}
\setcopyright{none}
\citestyle{acmnumeric}
\renewcommand{\bibliofont}{\normalfont\normalsize}
\newtheorem{theorem}{Theorem}
\newtheorem{lemma}[theorem]{Lemma}
\newtheorem{corollary}[theorem]{Corollary}
\theoremstyle{definition}
\newtheorem{definition}[theorem]{Definition}
\newcommand{\MST}{\operatorname{MST}}
\newcommand{\cost}{\operatorname{cost}}
\newcommand{\calC}{\mathcal C}
\newcommand{\calE}{\mathcal E}
\newcommand{\calG}{\mathcal G}
\newcommand{\core}{\operatorname{core}}
\newcommand{\ceil}[1]{\left\lceil #1\right\rceil}
\title[Deterministic MST Approximation under Decrease-Only Corruption]{Deterministic Quasipolynomial MST Approximation under Decrease-Only Corruption}
\author{Zijian Zeng}
\affiliation{\institution{UCSI University}\country{Malaysia}}
\email{1002266693@ucsiuniversity.edu.my}
\authorsaddresses{\normalsize Zijian Zeng, UCSI University,
\href{mailto:1002266693@ucsiuniversity.edu.my}{1002266693@ucsiuniversity.edu.my}}
\renewcommand{\shortauthors}{Zijian Zeng}
\date{}
\begin{document}
\begin{abstract}
An algorithm observes a finite metric after an adversary has decreased an
independently marked subset of its pairwise distances. The adversary knows
the metric and the entire mask, and the observations need not form a metric.
For every fixed marking probability $p<1$, we give a deterministic algorithm
that uses no true-distance queries and returns a spanning path of expected
true cost $O_p(\log^2 n)$ times the minimum spanning tree cost. On every input,
its running time is $n^{O_p(\log n)}$ in the exact arithmetic/comparison model,
independently of the true and observed aspect ratios.

The construction builds on the hidden-core hierarchy of Anand, Gupta,
Jayaram, Lee and Vakilian. Their quasipolynomial guarantee has approximation
factor $O(\log^4 n)$ for $p\leq0.1$; their unrestricted-time algorithm gives
$O_p(\log n)$ for all fixed $p<1$. A uniform uncorrupted-rectangle argument
and high quantiles extend hierarchical carving to every fixed positive
uncorrupted density. A deterministic depth-first leaf path admits an
outlier-sensitive charging bound: its cost depends on the actual weighted
outlier mass rather than a worst-case outlier allowance per cluster.
Separate packing bounds for generated clusters and leftover singletons
then give the quadratic logarithm.
\end{abstract}
\maketitle
\clearpage

\section{The computational question and the result}

Minimum spanning trees are easy to compute when their edge lengths are
known. Corrupted distances introduce a different obstacle: a very small
observed length can conceal a long true edge, and the observed matrix can
violate every useful triangle inequality. The strong/weak distance-oracle
framework of Bateni, Dharangutte, Jayaram and Wang~\cite{bateni2024}
motivates recovering a low-cost tree with few true-distance queries.
Their $O(\sqrt{\log n})$ MST upper bound assumes metric weak observations
(Theorems~4--6 distinguish this assumption from unrestricted observations).
Here all observations are available, but no true-distance queries are allowed.

Let $X=\{1,\ldots,n\}$, $n\geq2$, and let $d$ be a genuine finite metric:
$d(u,v)>0$ for distinct points. Independently mark each unordered pair with
probability $p$. An omniscient adversary sees $d$ and the whole mask, and
chooses finite symmetric observations $0\leq D_{uv}\leq d(u,v)$ on marked
pairs. Unmarked pairs have $D_{uv}=d(u,v)$; put $D_{uu}=0$.
Only $D$ is given to the algorithm. Write
$\cost_d(T)=\sum_{uv\in T}d(u,v)$ and $W=\MST(d)$.

Anand, Gupta, Jayaram, Lee and Vakilian~\cite{anand2026}, Theorems~1.2--1.3,
give an unrestricted-time $O_p(\log n)$ approximation for every fixed $p<1$
and an $n^{O(\log n)}$-time $O(\log^4 n)$ approximation for $p\leq0.1$.
The computational question is to obtain a polylogarithmic approximation
in quasipolynomial time throughout the decrease-only noise range.

\begin{theorem}\label{thm:main}
For every rational $0<q\leq1$ there are a deterministic algorithm and a
constant $a_q>0$ with the following properties. On every observed input
the algorithm returns a spanning path, uses no true-distance or mask-label
queries, and performs at most $n^{a_q\log_2 n}$ exact arithmetic/comparison
and control-bit operations. The bound is independent of aspect ratio.
For every metric $d$, every $0\leq p\leq1-q$, and every adversary described
above,
\[
 \mathbb E[\cost_d(P)]\leq 10^{15}q^{-5}(\log_2 n)^2\MST(d).
\]
The expectation is over the mask; the algorithm uses no private bits.
\end{theorem}

For each fixed real $0\leq p<1$, hardwiring a rational $0<q\leq1-p$
gives expected $O_p(\log^2 n)$ approximation in $n^{O_p(\log n)}$ operations.
Arithmetic on input distances is charged exactly; a bit-complexity claim
for arbitrary-precision rational distances would additionally charge their
precision.

\paragraph{Method and relation to the hierarchy framework.}
The hierarchy, universal future-seed consistency tests, normalization, and
promotion/intersection framework come from Section~6 of~\cite{anand2026}.
We retain this architecture. The changes address two losses within it.
First, even when most pairs are marked, every large observed clique has
an almost-full hidden core: one dense unmarked rectangle gives many true
neighbors of a center, and a second rectangle excludes many distant seed
points. High quantiles can then compare centers at any fixed unmarked
density. Universal future tests ensure that outliers do not accumulate
when clusters merge.

Second, we connect the leaves of the cluster tree in depth-first order.
Its final degree two controls delayed charges, and the contiguous interval
of each child permits at most one immediate right-exit charge. For weighted
cluster count $H$ and weighted actual outlier mass $\Omega$, the resulting
cost is at most $2\beta H+12\beta\Omega$ (Lemma~\ref{lem:extract}).
Both quantities have the same packing bound. In particular, a leftover
singleton contributes one outlier, even when generated clusters have a
logarithmic outlier allowance. These steps form one algorithmic method.

Related hierarchy-to-tree work by Fletcher and Hershkowitz~\cite{fletcher2026}
concerns graph-spanning distance embeddings; their stated constant-distortion
guarantee uses level-diameter growth $\Omega(\log n)$. Our objective is total
path cost in a metric closure, with dyadic levels and hidden cores that may
change nonmonotonically. This comparison uses the theorem stated in the
publisher abstract.

\section{Uniform mask control and seed quantiles}\label{sec:seeds}

Set $h=\log_2 n$, $\ell=\ceil{h}$, and
\[
 r=\ceil{80q^{-2}\ell},\qquad R=\ceil{128r/q},\qquad
 \tau=1-q/8,\qquad A=32,\qquad\beta=256.
\]
Then $r\leq81q^{-2}\ell$, $R\leq11000q^{-3}\ell$, and $r/R\leq q/128$.
The case $n=2$ uses its unique edge. If $R>n$, use the label-order path;
its cost is at most $(n-1)W<RW$, within Theorem~\ref{thm:main}.
For now assume $R\leq n$.

Let $\calE$ be the mask event that every two disjoint sets $U,V$ of size
at least $r$ have at least $(q/2)|U||V|$ unmarked cross pairs.
For fixed sizes $a,b$, Hoeffding's inequality bounds a rectangle's failure
probability by $\exp(-q^2ab/2)$. Since $ab\geq r(a+b)/2$,
\begin{equation}\label{eq:mask}
 \Pr(\calE^c)\leq
 \sum_{a,b=r}^n\exp\big((a+b)\ln n-q^2r(a+b)/4\big)
 \leq\sum_{a,b\geq1}n^{-27(a+b)}<n^{-20}.
\end{equation}
Indeed $q^2r/4\geq20\ell$ and $20/\ln2-1>27$.
When the disjoint sets do not exist, the condition is vacuous.
This event depends only on the mask. Its consequences therefore hold for
all adversarial observations and all adaptively selected seeds simultaneously.

The following lemmas apply to any metric and compatible decreased observations;
in the algorithm they will be applied after normalization.
A \emph{seed at scale $s$} is an $R$-set $S$ with $D_{uv}\leq s$
for all distinct $u,v\in S$.

\begin{lemma}[Seed core]\label{lem:seed}
On $\calE$, every seed $S$ at scale $s$ has a center $z_S\in S$ such
that $G_S=\{x\in S:d(x,z_S)\leq2s\}$ has size at least $R-r$.
\end{lemma}
\begin{proof}
Split $S$ into two halves, each of size at least $r$. The rectangle condition
gives a point $z$ in the first half with at least $qR/4\geq32r$ unmarked
neighbors $U$ in the second. All these neighbors are truly within $s$ of $z$.
If an $r$-set $F\subseteq S$ were farther than $2s$ from $z$, then $U,F$
would be disjoint and every true $U$--$F$ distance would exceed $s$.
Every such observation is at most $s$, so all these pairs would be marked,
contradicting $\calE$.
\end{proof}

For each $R$-set, fix a hidden center at the smallest scale at which it is
an observed seed. Its core-radius bound is then valid at every later
qualifying scale. These centers are analysis devices only.
For a list of $M$ numbers, its $\tau$-quantile means order statistic
$\ceil{\tau M}$. Define $Q(x,S)$ from the $R$ entries $D_{xs}$,
including the diagonal zero if $x\in S$. Define $Z(S,S')$ from all $R^2$
ordered entries $D_{ss'}$. Overlap can introduce diagonal zeros and repeated
unordered pairs; the definition is unchanged, including when $S=S'$.

\begin{lemma}[Quantile bounds]\label{lem:quantiles}
On $\calE$, seeds $S,S'$ with core radii $2s,2s'$ satisfy
\begin{align}
 Q(x,S)&\leq d(x,z_S)+2s,\label{eq:pointupper}\\
 |Z(S,S')-d(z_S,z_{S'})|&\leq2s+2s'.\label{eq:seedcompare}
\end{align}
Moreover, for any $B\geq0$, fewer than $r$ points $x\in X$ satisfy both
$Q(x,S)\leq B$ and $d(x,z_S)>B+2s$.
\end{lemma}
\begin{proof}
At least $R-r>\tau R$ point-to-core entries give~\eqref{eq:pointupper}.
At least $(R-r)^2>\tau R^2$ core-to-core entries give the upper half
of~\eqref{eq:seedcompare}. Both use $D\leq d$.
For its lower half put $b=d(z_S,z_{S'})-2s-2s'$. If $b\leq0$,
nonnegativity suffices. Otherwise the two cores are disjoint even if
the seeds overlap. The event $\calE$ gives at least
$(q/2)(R-r)^2>(1-\tau)R^2$ unmarked core pairs, each with observation
at least $b$. These are distinct unordered pairs. Thus fewer than
$\tau R^2$ entries are strictly below $b$, and the specified ceiling-rank
quantile is at least $b$. This also handles integer rounding.

For the final assertion suppose an $r$-set $F$ of such far points existed.
It is disjoint from $G_S$. For each $x\in F$, at most
$R-\ceil{\tau R}\leq qR/8$ observations exceed $B$. Every unmarked edge
from $x$ to $G_S$ must exceed $B$, by the triangle inequality. Hence the
rectangle has at most $|F|qR/8$ unmarked pairs. Its required number is at
least $(q/2)|F||G_S|>|F|qR/8$, since $|G_S|\geq R-r>R/4$.
\end{proof}

\section{The observable hierarchy}\label{sec:algorithm}

Compute $U=\max_{u<v}D_{uv}$. If $U=0$, return the label-order path.
Otherwise set
\[
 \lambda=2U/n^2,\qquad D^*_{uv}=\max\{1,D_{uv}/\lambda\}\quad(u\ne v),
 \qquad D^*_{uu}=0,\qquad L=\ceil{\log_2(n^2)}.
\]
All predicates below use $D^*$, with scales $t_i=2^i$. For analysis the
corresponding true metric is $d^*_{uv}=\max\{1,d(u,v)/\lambda\}$ off the
diagonal, and zero on it. Section~\ref{sec:expectation} verifies this
normalization without any aspect-ratio assumption.

For every seed $S$ at $t_i$, define $F_i(S,x)$ to mean that for every
$j\in\{i,\ldots,L\}$ and every $R$-set $S'$ that is a seed at $t_j$,
\begin{equation}\label{eq:future}
 Q(x,S')\leq At_i+Z(S,S')+4t_j.
\end{equation}
All seeds in the universe are tested, irrespective of their eventual selection
or intersection with $S$. A fresh seed $S$ is \emph{admissible} at $i$ if
it is a seed at $t_i$ and $F_i(S,x)$ holds for every $x\in S$.

Start with the singleton partition $\calC_0$. For each $i=1,\ldots,L-1$,
initialize an active list with all objects in $\calC_{i-1}$. Every
nonsingleton object carries an $R$-set representative; a singleton carries
its point. Build $\calC_i$ by the following deterministic loop.
\begin{enumerate}
\item If an active nonsingleton exists, select its lexicographically first
representative $S$. Otherwise select the lexicographically first admissible
$R$-set of active singleton points. If no such set exists, stop.
\item Form one output cluster from the selected old object, if any; all
other active old objects whose representative $S_0$ satisfies
$Z(S,S_0)\leq At_i$; and all active singletons $x$ satisfying
$Q(x,S)\leq At_i$ and $F_i(S,x)$.
\item Remove these entire objects from the active list. Give their union
representative $S$ and hold it separately until this level ends.
\end{enumerate}
At termination, $\calC_i$ consists of the held output clusters and leftover
singletons. Fresh seed members are included because $Q(x,S)\leq t_i$ and
admissibility supplies $F_i(S,x)$. Each iteration removes an object; outputs
never reenter the same level's list. Consequently these are nested partitions,
and every nonsingleton contains its representative seed, which remains a
clique at larger scales. Set $\calC_L=\{X\}$.

Treat each level-cluster pair as a node, including repeated underlying sets.
Order children by their least label. List the singleton leaves in depth-first
order and connect consecutive leaves. This specifies a spanning path on every
input using observed data alone.

\paragraph{A useful completeness implication.}
For the rest of Sections~\ref{sec:algorithm}--\ref{sec:packing}, distances
mean $d^*$ and quantiles mean those of $D^*$. On $\calE$, if
$d^*(x,z_S)\leq(A-2)t$ for a seed at current scale $t$, then $x$ passes
both singleton tests. Equation~\eqref{eq:pointupper} gives $Q(x,S)\leq At$.
For every future seed at scale $t'$,
\begin{align*}
 Q(x,S')&\leq d^*(x,z_S)+d^*(z_S,z_{S'})+2t'\\
 &\leq(A-2)t+Z(S,S')+2t+4t'=At+Z(S,S')+4t'.
\end{align*}
The current radius $2t$ is valid for carried seeds as well.
In particular any $R$-set of true diameter at most $t$ is admissible.

\begin{lemma}[Nonaccumulating outliers]\label{lem:nonaccum}
On $\calE$, every non-root nonsingleton $C\in\calC_i$ has a hidden
core of diameter at most $256t_i$, size greater than $r$, and fewer
than $r$ outliers. Singletons may be assigned empty cores, with one outlier.
\end{lemma}
\begin{proof}
Let $S_f$ represent $C$ at scale $t$. Each $x\in C$ first ceased to be
a singleton at a scale $t_0\leq t$, with selected representative $S_0$.
The singleton test (also for fresh seed members) gave $F_{i_0}(S_0,x)$.
Follow the representatives up the ancestor chain to $C$. There is at most
one change per level. At a change at scale $s$, the merge rule and
\eqref{eq:seedcompare} bound the center displacement by $(A+4)s$.
Thus
\[
 d^*(z_{S_0},z_{S_f})\leq(A+4)\sum_{\text{dyadic }s\leq t}s
 \leq2(A+4)t.
\]
The final seed $S_f$ was among the universal future tests at scale $t$,
even before it was selected. Its fixed birth center has radius at most $2t$;
the birth center of $S_0$ has radius at most $2t_0$. Therefore
\begin{align*}
 Q(x,S_f)&\leq At_0+Z(S_0,S_f)+4t\\
 &\leq At+2(A+4)t+2t_0+2t+4t\leq112t.
\end{align*}
The setwise soundness in Lemma~\ref{lem:quantiles} says that fewer than
$r$ points in the entire set $\{x:Q(x,S_f)\leq112t\}$ lie farther than
$114t$ from $z_{S_f}$. Removing these points from $C$ leaves diameter
at most $228t<256t$ and size at least $R-r>r$.
The bound concerns the whole merged cluster and does not add incoming
outlier allowances.
\end{proof}

\section{Packing with the actual outlier mass}\label{sec:packing}

Write $W^*=\MST(d^*)$. Selected representatives at one scale $t$ have
centers more than $(A-4)t=28t$ apart. To see this, suppose an earlier
selection $S$ and later selection $S'$ had closer centers. If $S'$ was an
active old object, \eqref{eq:seedcompare} would give $Z(S,S')\leq At$,
so that object would have been removed. If it was fresh, its nonempty
seed core consists of singleton points within $2t+28t=(A-2)t$ of $z_S$.
The completeness implication would have removed these points, contradicting
their later availability.

In any metric, $k$ points separated by more than $\delta$ satisfy
$k\leq1+2\MST(d)/\delta$. For $k\geq2$, shortcutting a doubled MST
gives a cyclic tour through the points of length at most $2\MST(d)$,
whereas each of its $k$ steps exceeds $\delta$; $k=1$ is immediate.
Hence the number $K_i$ of seed-generated clusters is at most
$1+W^*/t_i$.

No true-diameter-$t_i$ set contains $R$ leftover singletons, since those
points would be an admissible seed and the loop could not stop. Cover
leftover points greedily by true radius-$t_i/2$ balls centered at uncovered
leftovers. The centers are more than $t_i/2$ apart, so at most
$1+4W^*/t_i$ balls are used, each containing at most $R-1$ leftovers.
If there are none, use zero balls. Writing $N_i$ for the number of leftover
singletons gives
\begin{equation}\label{eq:counts}
 K_i\leq1+W^*/t_i,\quad N_i\leq(R-1)(1+4W^*/t_i),\quad
 |\calC_i|\leq R+4RW^*/t_i.
\end{equation}
These two categories exhaust the partition, including carried old clusters:
every old nonsingleton is selected or absorbed before the loop stops.

Assign each intermediate nonsingleton its core from
Lemma~\ref{lem:nonaccum}, each singleton the empty core, and the root the
full core $X$. Put
\[
 H=\sum_{i=0}^L t_i|\calC_i|,\qquad
 O_i=\sum_{C\in\calC_i}|C\setminus\core_i(C)|,\qquad
 \Omega=\sum_{i=0}^L t_i O_i.
\]
For intermediate levels, \eqref{eq:counts} gives the sharper category bound
\begin{equation}\label{eq:outliers}
 O_i\leq rK_i+N_i\leq2R+5RW^*/t_i.
\end{equation}
At the bottom $O_0=n$; at the root $O_L=0$.

Assume temporarily that the true normalized minimum distance is at least
$1$ and its diameter $\Delta^*$ belongs to $[n^2/2,n^2]$. Then
$W^*\geq n-1$, $n\leq2W^*$, and
$t_L<2n^2\leq4\Delta^*\leq4W^*$. Summing
\eqref{eq:counts} and \eqref{eq:outliers}, using
$\sum_{i=1}^{L-1}t_i<t_L$, yields
\begin{align}
 H&\leq n+(R+1)t_L+4R(L-1)W^*
       \leq14R(L+1)W^*,\label{eq:H}\\
 \Omega&\leq n+2Rt_L+5R(L-1)W^*
       \leq14R(L+1)W^*.\label{eq:Omega}
\end{align}
The singleton contribution in \eqref{eq:outliers} is one per point.
Replacing it by the maximum allowance $r$ would lose a logarithmic factor.

\section{Deterministic extraction from hidden cores}\label{sec:extract}

The following statement isolates the extraction step from the carving
algorithm. Cores need neither be computable nor be nested.

\begin{lemma}[Outlier-sensitive leaf path]\label{lem:extract}
Let $\calC_0,\ldots,\calC_L$ be nested partitions of a finite metric,
with singleton bottom and one-cluster top. At scale $t_i=2^i$, each core
has diameter at most $\beta t_i$. The root core is $X$. Every non-root
nonempty core has size greater than $b$ and omits at most $b$ points of
its cluster, where $b\geq1$; each empty-core cluster has at most $B$ points,
where $B\geq b$. Define $H,\Omega$ as above. Every compatible depth-first
ordering of the singleton leaves induces a path $P$ with
\begin{equation}\label{eq:extraction}
 \cost(P)\leq2\beta H+12\beta\Omega.
\end{equation}
\end{lemma}
\begin{proof}
Orient each path edge from its earlier leaf $u$ to its later leaf $v$.
Let $j\geq1$ be their least common cluster level. Promote each endpoint
$x$ from level $j$ to the first containing cluster $A_x$ whose core
contains $x$, at level $k_x$. These levels exist because the root core
is $X$. Put $k=\max\{k_u,k_v\}$ and $D=\calC_k(u)=\calC_k(v)$,
where $\calC_i(x)$ denotes the cluster containing $x$.
The core of $D$ is nonempty since the later-promoted endpoint belongs to it.

The cores of $A_x$ and $D$ intersect. This is immediate if $A_x=D$ or
if $D$ is the root. Otherwise $A_x$ is non-root, its core has more than
$b$ points, and it is contained in $D$, which omits at most $b$ points.
Choose $w_x$ in this intersection. Three diameter bounds give
\[
 d(u,v)\leq d(u,w_u)+d(w_u,w_v)+d(w_v,v)
 \leq\beta2^{k_u}+\beta2^k+\beta2^{k_v}\leq3\beta2^k.
\]
If $k=j$, both endpoints lie in the common level-$j$ core, giving the
better bound $d(u,v)\leq\beta2^j$. Charge the edge to the level-$(j-1)$
child containing $u$. This child occupies a contiguous interval in the
leaf order; an edge leaves its right boundary at most once. These charges
sum to at most $2\beta H$.

If $k>j$, choose an endpoint $x$ with $k_x=k$, breaking a tie by $u$.
It is an outlier in $\calC_{k-1}(x)$ by minimality of its promotion level.
Charge that cluster; the edge costs at most $6\beta2^{k-1}$. Each outlier
point has at most two incidences in the final path. Thus a level-cluster
pair receives at most twice its actual outlier count in charges, for a
total at most $12\beta\Omega$. The two edge classes prove the bound.
Repeated underlying clusters are distinct level nodes; no step assumes
that core membership persists to the next level.
\end{proof}

For the constructed hierarchy use $b=B=r$ and $\beta=256$.
Lemma~\ref{lem:nonaccum} supplies all intermediate core conditions.
Under the temporary scale assumptions the root has diameter at most
$t_L$, so \eqref{eq:H}--\eqref{eq:extraction} imply
\begin{equation}\label{eq:normalized}
 \cost_{d^*}(P)\leq196\beta R(L+1)W^*
                  =50176R(L+1)W^*.
\end{equation}
The path is the leaf order of the cluster tree; no hidden core is needed
to compute it.

\section{Aspect normalization and unconditional expectation}\label{sec:expectation}

Let $\Delta$ be the original diameter. Fix a diametral pair $a,b$.
For every other $x$, choose one of $ax,bx$ of length at least $\Delta/2$.
Together with $ab$, these are $n-1$ distinct pairs depending only on $d$.
If any is unmarked, then $U\geq\Delta/2$. Thus the scale event
$\calG=\{U\geq\Delta/2\}$ satisfies
\begin{equation}\label{eq:scale}
 \Pr(\calG^c)\leq p^{n-1}
\end{equation}
uniformly over the adversary. This also includes the $U=0$ branch.

For $U>0$, the normalized $d^*$ is a genuine metric. For three distinct
points, if a left-hand distance is $1$, the other two are at least $1$;
otherwise the original triangle inequality and monotonicity suffice.
Repeated-vertex triangle inequalities are identities. Also $D^*\leq d^*$,
and unmarked equality persists. The same mask event $\calE$ therefore
controls the transformed instance even though $\lambda$ depends on $D$.
On $\calG$, since $U\leq\Delta$,
\[
 n^2/2\leq\Delta^*=\max\{1,\Delta/\lambda\}\leq n^2.
\]
This establishes the scale assumptions used above. Moreover
\[
 W^*\leq W/\lambda+(n-1),\qquad
 \lambda(n-1)\leq2\Delta/n\leq2W/n,\qquad
 d(u,v)\leq\lambda d^*(u,v).
\]
For the first inequality evaluate an original MST in $d^*$; for the
second use $\lambda=2U/n^2$ and $U\leq\Delta$.
Consequently \eqref{eq:normalized} gives, on $\calE\cap\calG$,
\[
 \cost_d(P)\leq100352R(L+1)W.
\]
On every other mask the output is still a spanning path, with cost
at most $(n-1)\Delta<nW$. By \eqref{eq:mask} and \eqref{eq:scale},
their unconditional expected contribution is at most
\[
 n(n^{-20}+p^{n-1})W\leq(1+1/q)W.
\]
Indeed $np^{n-1}\leq\sum_{k=0}^{n-1}p^k\leq1/(1-p)\leq1/q$;
this includes $p=0$ literally for $n\geq2$. No conditional distribution
of the mask is used.

Finally $L\leq2\ell$, $L+1\leq3\ell$, and $\ell\leq2h$.
The good-event term is bounded by
\[
 100352\cdot11000\cdot3\cdot4\,q^{-3}h^2 W
   =13246464000\,q^{-3}h^2 W.
\]
Adding $(1+1/q)W$ is within $10^{15}q^{-5}h^2W$.
The $n=2$ and $R>n$ branches already satisfy that bound, completing
the approximation part of Theorem~\ref{thm:main}.

\section{Every-execution operation bound}\label{sec:runtime}

Integer doubling computes $\ell,L$; fixed rational arithmetic computes
$r,R$ and all quantile ranks. No real logarithm or rounding oracle is needed.
In the nontrivial branch put $N=\binom nR\leq n^R$ and enumerate all
$R$-subsets. Check their clique status at every scale. Precompute $Q(x,S)$
and $Z(S,S')$ by sorting lists of lengths $R$ and $R^2$ and selecting
the exact ceiling ranks. Even quadratic sorting takes $O(R^4)$ comparisons
per seed pair. Scanning all future scales and all clique seeds computes
the $F_i(S,x)$ tables with $O(L^2nN^2)$ further comparisons. These tables
include identical and intersecting seeds.

There are at most $n$ carving iterations per level. Each fresh-seed search
can scan all $N$ subsets and their active memberships. Label lists or bit
vectors represent objects, their unions and the partition parents; copying,
tie breaking, leaf traversal and writing the path have polynomial overhead.
Since $R\leq n$ and $L\leq2\ell\leq2n$, a fixed implementation uses at
most $cL^2n^{2R+6}$ elementary operations before charging integer-control
overhead, for an absolute constant $c$.

Control variables have $O(R\log n+\log L)$ bits plus a $q$-dependent
constant. Charging their elementary bit operations adds polynomial factors
in $R,\log n$, so fixed routines have a bound $c_q n^{2R+k}$ for constants
$c_q,k$. As $R\leq22000q^{-3}h$, choosing an integer
\[
 a_q\geq44000q^{-3}+k+\ceil{\log_2\max\{1,c_q\}}+1
\]
gives $c_qn^{2R+k}\leq n^{a_qh}$ for $n\geq2$.
The constants are determined by the fixed enumeration, sorting, and integer
routines, independently of distances and their precision. All operations
on distance values are charged in the exact arithmetic/comparison model.
The loops depend only on $n,q$ and have the same bounds on bad masks.
There are no rejection loops or private random bits. The early branches
and the forced root ensure a valid output on every execution. This proves
the remaining assertions of Theorem~\ref{thm:main}.

\section{Discussion and verification}

The theorem combines the full fixed decrease-only noise range with a
quasipolynomial deterministic algorithm and a quadratic logarithmic
approximation. The unrestricted-time logarithmic benchmark remains the
nearest approximation target. Whether the same computational model admits
a logarithmic approximation, or whether the extraction inequality yields
other applications, are further questions.

The proof is analytic for all $n$. As supplementary checks, exact integer
shortest-path closures were used to test the hidden-core extraction argument
on finite hierarchies. A separately implemented check enumerated all 358
nested two-intermediate-level partition chains on five labels, with 8,951
eligible core assignments for $(b,B)=(1,2),(2,3)$ and 268,560 compatible
ordered cases. It checked 1,074,240 path edges, the promotion/intersection
conditions, and both charge bounds. These finite checks support debugging;
they do not replace Lemma~\ref{lem:extract} or execute the large-constant
carving algorithm.

\paragraph{AI methodology.}
OpenAI Codex was used substantially for algorithm and proof development,
exact finite-check code, literature comparison, and manuscript drafting.
Separate model contexts performed internal mathematical checks.

\bibliographystyle{ACM-Reference-Format}
\bibliography{references}
\end{document}
